\section{Machine-learning concepts used in the project}\label{app:theory}

\paragraph{A.1 Bias--variance trade-off.} For squared loss and $y=f(x)+\varepsilon$, the expected test
error of an estimator $\hat f$ trained on a random sample decomposes as
\[\E\big[(y-\hat f(x))^2\big]=\underbrace{\big(f(x)-\E\hat f(x)\big)^2}_{\text{bias}^2}+
\underbrace{\mathrm{Var}\big(\hat f(x)\big)}_{\text{variance}}+\sigma_\varepsilon^2 .\]
Flexible models (deep trees) have low bias and high variance; rigid models (penalised linear) the
opposite. The noise term $\sigma^2_\varepsilon$ is irreducible. In commodity markets it is very large
relative to the predictable part, so variance control (regularisation, averaging, small grids) matters more
than flexibility \citep[ch.~7]{hastie2009}.

\paragraph{A.2 Regularisation.} Ridge and $L_2$-logistic regression add $\lambda\lVert\beta\rVert_2^2$ to
the loss. In the linear case the solution $\hat\beta=(X^\top X+\lambda I)^{-1}X^\top y$ shrinks the
coefficients along the principal components of $X$, most strongly along low-variance directions, which
stabilises estimates when features are collinear. The penalty is not scale-invariant, so features are
standardised, using means and variances of the training window only.

\paragraph{A.3 Decision trees, bagging and random forests.} A CART tree chooses, at each node, the
feature and threshold that most reduce impurity (Gini or entropy for classification, squared error for
regression) and predicts the average outcome in each leaf. Trees are invariant to monotone transformations
of features and model interactions automatically, but small changes in the data change the tree a lot.
Bagging \citep{breiman1996} averages trees fitted on bootstrap samples; if each tree has variance $\sigma^2$
and pairwise correlation $\rho$, the average of $B$ trees has variance
$\rho\sigma^2+(1-\rho)\sigma^2/B$. As $B$ grows only $\rho\sigma^2$ remains, so random forests also
sample a subset of features at each split to decorrelate the trees \citep{breiman2001}. Because our labels
overlap in time, bootstrap samples contain near-duplicate rows; we subsample 50\% of rows per tree and use
large minimum leaf sizes to limit this (a ``sequential bootstrap'' as in \citealp{lopezdeprado2018} would
be a refinement).

\paragraph{A.4 Gradient boosting.} Boosting minimises $\sum_t L(y_t,F(x_t))$ by functional gradient
descent: starting from a constant $F_0$, at step $m$ compute pseudo-residuals
$r_{tm}=-\partial L(y_t,F)/\partial F|_{F=F_{m-1}(x_t)}$, fit a small regression tree $h_m$ to them, and update
$F_m=F_{m-1}+\nu h_m$ \citep{friedman2001}. For the log loss on the log-odds scale, $r_{tm}=y_t-\hat\pi_{m-1}(x_t)$.
A small learning rate $\nu$ with more trees acts as regularisation (shrinkage); the depth of the trees
bounds the order of interactions.

\paragraph{A.5 Proper scoring rules and probability evaluation.} A scoring rule $S(p,y)$ is strictly
proper if $\E_{y\sim q}S(p,y)$ is uniquely optimised at $p=q$ \citep{gneiting2007}. The log loss and the
Brier score $(p-y)^2$ are both strictly proper, so they reward honest probabilities; accuracy is not.
The Brier skill score $1-\mathrm{BS}/\mathrm{BS}_{\text{ref}}$ compares a model with the climatological
forecast (the training base rate): positive values mean the model adds information. The ROC AUC equals the
probability that a randomly chosen spike day receives a higher score than a randomly chosen non-spike day
\citep{fawcett2006}; it measures ranking (discrimination) and ignores calibration, which is shown
separately with a reliability diagram.

\paragraph{A.6 Cross-validation for overlapping financial labels.} Standard $K$-fold CV assumes
independent observations. With time series, a random split puts tomorrow's data in the training set and
yesterday's in the validation set; with \emph{overlapping} labels (each label uses the next 22 days) even a
chronological split leaks, because the last training labels are computed from prices inside the validation
period. \emph{Purging} removes every training row whose label window ends after the validation period
starts; an \emph{embargo} removes a further block of rows, because features are also autocorrelated
\citep[ch.~7]{lopezdeprado2018}. We apply both in the outer walk-forward loop and in the inner
cross-validation used to tune hyper-parameters, and test them automatically (\texttt{tests/test\_leakage.py}).

\paragraph{A.7 Permutation importance.} The importance of feature $j$ is the drop in out-of-sample
performance (here AUC) when column $j$ of the test set is randomly permuted, which breaks its link with the
target while keeping its distribution \citep{fisher2019}. Unlike impurity-based importance it is computed
on test data and is not biased towards high-cardinality features \citep{strobl2007}. With correlated
features it understates the importance of each member of a correlated group, so we also report importance
summed by group.

\paragraph{A.8 Diebold--Mariano test.} To compare two forecasts with errors $e_{1t},e_{2t}$, define
$d_t=e_{1t}^2-e_{2t}^2$. Under the null of equal accuracy $\E d_t=0$, and
$\mathrm{DM}=\bar d/\sqrt{\widehat{\mathrm{LRV}}(d)/n}\to\mathcal N(0,1)$ \citep{diebold1995}. Daily forecasts of a
21-day average are autocorrelated by construction, so the long-run variance uses a Newey--West estimator
with 21 lags \citep{newey1987}.

\paragraph{A.9 Conditional Value-at-Risk.} For a loss $X$, $\mathrm{CVaR}_{95}(X)=\E[X\mid X\ge
\mathrm{VaR}_{95}(X)]$, the average of the worst 5\% outcomes \citep{rockafellar2000}. Unlike variance, it
only penalises bad outcomes, which matches the haulier's concern: months when diesel costs far more than
the price that could have been locked.

\paragraph{A.10 Realised volatility and the Garman--Klass estimator.} Volatility is not observed; it is
estimated from prices. Squared daily returns are unbiased but very noisy, and with generic futures series
they are contaminated by contract rolls. Range-based estimators use the open $O$, high $H$, low $L$ and close
$C$ of the same contract on the same day. \citet{garman1980} show that
$\hat\sigma^2_{GK}=\tfrac12\ln^2(H/L)-(2\ln2-1)\ln^2(C/O)$ is unbiased under a driftless random walk and up to
seven times more efficient than the squared close-to-close return. Averaging it over 21 days gives our
realised-volatility target. It ignores overnight moves, which are small for futures that trade almost 24
hours, larger for agricultural contracts.

\paragraph{A.11 The HAR model.} \citet{corsi2009} models tomorrow's (log) realised volatility as a linear
function of its daily, weekly and monthly averages,
$\ln\sigma_{t+1:t+h}=\beta_0+\beta_d\ln\sigma_t^{(1)}+\beta_w\ln\sigma_t^{(5)}+\beta_m\ln\sigma_t^{(22)}+\varepsilon$.
The cascade reproduces the long memory of volatility with three parameters, and it is the standard
benchmark that machine-learning volatility models must beat \citep{christensen2023}. Adding the log implied
volatility (HAR-IV) brings in the market's forward-looking forecast; implied volatility exceeds
subsequent realised volatility on average (the variance risk premium) but contains information not in past
returns \citep{bollerslev2009}.

\paragraph{A.12 Neural networks.} The multilayer perceptron computes
$f(x)=W_3\,\phi(W_2\,\phi(W_1x+b_1)+b_2)+b_3$ with ReLU activations $\phi(z)=\max(0,z)$; with enough units it
can approximate any continuous function. It is trained by stochastic gradient descent (Adam) on the squared
loss with an $L_2$ penalty $\alpha$ on the weights, and with \emph{early stopping}: 10\% of the training rows
are held out and training stops when their loss no longer improves, a form of regularisation. Inputs must be
standardised. Neural networks need many observations; pooling 22 markets is what makes them usable here
\citep{gu2020}.

\paragraph{A.13 Lasso.} The Lasso adds $\alpha\lVert\beta\rVert_1$ to the squared loss. Unlike ridge, the
$L_1$ penalty sets some coefficients exactly to zero (the solution sits on a corner of the constraint set), so
it performs variable selection; with correlated features it tends to keep one member of a group.

\paragraph{A.14 Forecast combination.} Averaging forecasts of different models reduces variance when their
errors are imperfectly correlated, by the same argument as bagging (A.3). Equal weights are hard to beat
because estimated weights add noise; our ``ML ensemble'' is the simple mean of the four ML forecasts, fixed
in advance.

\paragraph{A.15 Loss functions for volatility and coverage tests.} Realised volatility is a noisy proxy
of true volatility. \citet{patton2011} shows that only some loss functions rank forecasts consistently when
the proxy is noisy: the squared error and the QLIKE loss $\hat\sigma^{-2}\sigma^2-\ln(\hat\sigma^{-2}\sigma^2)-1$,
which penalises under-prediction of risk more than over-prediction. For the 95\% cost bound we use the
\citet{kupiec1995} test: with $x$ breaches in $N$ months, $\mathrm{LR}=-2\ln\frac{0.95^{N-x}0.05^{x}}
{(1-\hat p)^{N-x}\hat p^{x}}\sim\chi^2_1$ under correct coverage, $\hat p=x/N$.
